
In view of \SGARefLink{sga3:xvii:lemma:3-1:item:ii}{3.1(ii)} and (iii), we may suppose that \(k\) is
algebraically closed and that \(V\) is finite-dimensional over \(k\).

Let
\[
  0\longrightarrow G'\longrightarrow G\longrightarrow G''
    \longrightarrow0
\]
be an exact sequence of algebraic groups, and suppose that \(G\) acts
on a sheaf \(V\) (for the fpqc topology).  Of course
\(V^G\subset V^{G'}\), and the quotient presheaf \(G/G'\) acts
naturally on \(V^{G'}\).  But \(V^{G'}\) is a sheaf, being the kernel
of the familiar pair of morphisms
\[
  V\rightrightarrows\underline{\operatorname{Hom}}(G',V).
\]
Consequently the sheaf associated with \(G/G'\), namely \(G''\),
acts on \(V^{G'}\), and it is immediate that
\[
  \left(V^{G'}\right)^{G''}
  =V^G
  =\left(V^{G'}\right)^{G/G'}.
\]

This remark reduces the proof of \SGARefLink{sga3:xvii:proposition:3:2}{3.2} to the case in which \(G\) is an
elementary unipotent group (\SGARefLink{sga3:xvii:corollary:1:7}{1.7}).

\SGARefTarget{sga3:xvii:proposition:3:2:case:a}
\textit{\SGARefLink{sga3:xvii:appendix:A}{(a)} \(G=\mathbb G_a\), \(p=0\).}
It follows from (BIBLE 4, Proposition 4) that a morphism from
\(\mathbb G_a\) to the linear group \(\operatorname{GL}(V)\) is given
by an exponential map
\[
  T\longmapsto\sum_{q=0}^{\infty}T^q\,\frac{n^q}{q!},
\]
where \(n\) is a nilpotent endomorphism of \(V\).  But then
\(V\ne0\) implies \(\operatorname{Ker}n\ne0\), and it is clear that
every vector of \(V\) annihilated by \(n\) is fixed by \(G\).

Suppose now that \(p>0\).

\SGARefTarget{sga3:xvii:proposition:3:2:case:b}
\textit{\SGARefLink{sga3:xvii:section:b}{(b)} \(G=\boldsymbol\alpha_p\).}
Since \(\boldsymbol\alpha_p\) is a radicial group of height \(1\)
(Exposé VII\textsubscript{A}, \S7), giving a representation of
\(\boldsymbol\alpha_p\) in \(V\) amounts to giving a representation
of the \(p\)-Lie algebra \(\boldsymbol\alpha_p\) in
\(\mathfrak{gl}(V)\) (Appendix II, \SGARefLink{sga3:ii:subsection:2-2}{2.2}), that is, in this case, an
element \(X\in\operatorname{End}(V)\) such that \(X^p=0\)
(Appendix II, \SGARefLink{sga3:xvii:proposition:2:1}{2.1}).  But then
\[
  V\ne0\quad\Longrightarrow\quad
  W=\operatorname{Ker}(X)\ne0,
\]
and, again by Appendix II, \SGARefLink{sga3:ii:subsection:2-2}{2.2}, one has
\(W=V^{\boldsymbol\alpha_p}\).

\SGARefTarget{sga3:xvii:proposition:3:2:case:c}
\textit{\SGARefLink{sga3:xvii:appendix:C}{(c)} \(G=\mathbb Z/p\mathbb Z\).}
A representation of \(G\) in \(V\) is equivalent to the choice of an
element \(x\in\operatorname{Aut}(V)\) such that \(x^p=1\), that is,
\((1-x)^p=0\).  Thus \(x\) is of the form \(1+n\), with \(n\)
nilpotent, and \(W=\operatorname{Ker}n\) is fixed by \(x\).

\SGARefTarget{sga3:xvii:proposition:3:2:case:d}
\textit{(d) \(G=\mathbb G_a\).}
Let \(G_i\), \(i\in I\), be the increasing filtered family of étale
algebraic subgroups of \(\mathbb G_a\), hence groups isomorphic to
\((\mathbb Z/p\mathbb Z)^{r_i}\) (\SGARefLink{sga3:xvii:proposition:1:5}{Proposition 1.5}), and put
\(V_i=V^{G_i}\).  Since \(V\) has finite nonzero dimension and each
\(V_i\) is nonzero by \SGARefLink{sga3:xvii:appendix:C}{(c)}, the decreasing filtered family of the
\(V_i\) is stationary, and
\[
  W=\bigcap_{i\in I}V_i\ne0.
\]
We shall use the following lemma.

\SGAStatement{Lemma 3.3}{The family of étale subgroups of
\((\mathbb G_a)_S\), where \(S\) is a prescheme of characteristic
\(p>0\), is schematically dense in \(G\) (\SGARefLink{sga3:ix:definition:4:1}{Exposé IX, 4.1}).}
\SGARefTarget{sga3:xvii:lemma:3:3}

By \SGARefLink{sga3:ix:lemma:4:4}{Exposé IX, 4.4}, it suffices to prove the lemma when \(S\) is the
spectrum of the prime field \(\mathbb F_p\).  In that case it is
enough to consider the family of étale subgroups \(G_n\), \(n\geq1\),
defined by
\[
  X^{p^n}-X=0.
\]
This family is schematically dense in
\((\mathbb G_a)_{\mathbb F_p}\), since it contains every closed point.

With this established, let us return to the proof of \SGARefLink{sga3:xvii:proposition:3:2:case:d}{3.2(d)}.  If
\(w\in W\), its stabilizer in \(G\) is an algebraic subgroup of
\(\mathbb G_a\) containing \(G_i\) for every \(i\in I\); it is
therefore equal to \(\mathbb G_a\) by \SGARefLink{sga3:xvii:lemma:3:3}{3.3}.  Consequently
\(W=V^G\).

\SGARefLink{sga3:xvii:proposition:3:2}{Proposition 3.2} immediately gives the following.
